% =====================================================================
% SECTION 6 - EVALUATION (DRAFT)
%
% Every number below is measured or derived, with its source noted in a
% comment. Two placeholders are marked PENDING and must be filled before
% submission. Nothing else is a placeholder.
% =====================================================================

\section{Evaluation}
\label{sec:evaluation}

We evaluate BVI against five questions: the on-chain cost of the anchor
(RQ1), the separability of the content digest (RQ2), the behaviour of path
attestation under partial carrier participation (RQ3), attack coverage (RQ4),
and verification latency (RQ5). Contract measurements were taken on a local
Hardhat chain with \texttt{solc}~0.8.24, the optimiser enabled at 200~runs, and
EVM version \texttt{paris}; each gas figure is the median of five runs of
\texttt{gasUsed} taken from the transaction receipt.

\subsection{On-chain cost (RQ1)}

Table~\ref{tab:gas} reports the cost of every registry operation. The figure
that matters is \texttt{anchorSession} at $91{,}535$~gas, since it is the only
per-session cost; Layer~1 operations are administrative and amortised over an
organisation's call volume, and Layer~2 attestation is optional and paid by the
participating carrier.

The anchor's cost is invariant in both call duration and hop count. Across a
grid of thirty duration and hop-count combinations spanning $30$~seconds to
$50$~minutes and one to twenty hops, measured cost varied by at most $36$~gas,
or $0.04$~percent --- consistent with block-state variation on a local chain
rather than with any dependence on the varied parameters. The mechanism is that
both arguments are fixed-width: a Merkle root over the digest sequence is
thirty-two bytes whatever the sequence length, and a hash over the attestation
set is thirty-two bytes whatever the set size.

% Source: results/table1_gas.json, results/b1_anchor_vs_bbca.json
Table~\ref{tab:anchorcost} sets this against BBCA. Counting writes from their
published Call Flow Control Table structure and charging each at BVI's own
per-write cost, BBCA's per-session overhead is $5\times$ at two hops and
$23\times$ at eight. The assumption favours BBCA, so these are lower bounds.

\subsection{Content digest separability (RQ2)}

% PENDING: regenerate on a real narrowband corpus before submission.
% Current figures are from synthetic source-filter signals and validate the
% pipeline only. See docs/FINDINGS.md.
Table~\ref{tab:digest} reports bit error rate between the digest computed at
the caller and the digest recomputed by the recipient, under benign channel
conditions and under media substitution. At the operating threshold --- set at
the $99$th percentile of benign scores, which is all a deployed verifier can
observe without seeing an attack --- wholesale substitution, replay at offset,
and time-reversal are detected without error.

Detection of partial splicing degrades with the substituted fraction: a splice
covering half a call is detected in $86$~percent of trials, a quarter in
$61$~percent, and five percent of a call in only $8$~percent. This is a real
limitation and we state it plainly. A short splice leaves most of the digest
matching, and the detection floor is set by the worst benign channel condition
the verifier must tolerate: admitting calls at $10$~dB~SNR raises the threshold
to $0.283$ and costs splice sensitivity, while requiring $20$~dB lowers it to
$0.220$ and recovers some.

The parameterisation is $17$~Bark-spaced bands over $300$--$3400$~Hz with a
$32$~ms frame and $20$~ms hop, giving $16$~bits per frame and a rate of
$800$~bps. That fits every side channel we considered, including a four-byte
RFC~8285 header extension at $1{,}600$~bps. Capacity is therefore not the
binding constraint; survivability is, since header extensions are stripped by
transcoding gateways. We treat this as a deployment question for the testbed
rather than a design question.

One finding from this experiment changed the design. Frames carrying no speech
produce a digest determined by rounding rather than by content, so any
perturbation randomises them; including them in the comparison drove benign
G.711 to a bit error rate of $0.31$, close to chance. Restricting comparison to
frames carrying speech reduces this to $0.012$. The consequence is that when a
window holds too little speech, the correct output is not a pass or a failure
but no verdict --- making Layer~3 ternary in precisely the sense Layer~2
already is, and for the same reason. Silence is not evidence of substitution.

\subsection{Path attestation under partial participation (RQ3)}

% Source: results/b2_participation_sweep.json, results/table3_path_coverage.json
Table~\ref{tab:participation} reports the distribution of path status over
$1{,}000$~simulated sessions at each participation level, with a path length of
five hops and thirty percent of sessions carrying an in-path identifier
rewrite.

Two results stand out. First, BVI yields usable evidence for every session at
every participation level, including zero, because the anchor commits Layers~1
and~3 irrespective of whether any carrier attested. BBCA yields a complete
record only when every hop is registered and participating: $4.9$~percent of
sessions at fifty percent participation, and none at zero.

Second, and contrary to the intuition that long international routes are the
hard case, BVI's detection \emph{improves} with path length while BBCA's
degrades. Reaching fifty percent detection of an in-path rewrite requires sixty
percent carrier participation on a three-hop path but only thirty percent on a
twelve-hop path; the corresponding BBCA figures are eighty percent and, at
twelve hops, unreachable. The mechanism is straightforward once stated: BVI
requires one participating hop on either side of the rewrite, and a longer path
offers more opportunities for that, whereas BBCA requires an unbroken chain and
a longer path offers more opportunities to break it.

The cost of this is a false demotion rate that plateaus near nine percent of
benign sessions, arising entirely from legitimate identifier rewrites performed
by unregistered session border controllers. This rate is flat in carrier
participation --- recruiting more carriers does not reduce it. The only lever
is SBC registration coverage, which is an operational finding rather than a
protocol one.

\subsection{Attack coverage (RQ4)}

% Source: test/scenarios.test.js, all rows asserted against the deployed contract
Table~\ref{tab:coverage} states coverage across ten attack classes for BBCA,
the digest alone, the detector alone, and BVI composed. Every row BVI claims to
block is asserted against the deployed contract in the accompanying test suite.

Two rows motivate the composition. An in-path identifier rewrite that leaves
the audio untouched defeats the digest and the detector completely, and is
caught only by path attestation. A media substitution in transit that leaves
the route untouched defeats path attestation, and --- if the substituted audio
is genuine human speech rather than synthesis --- defeats the detector as well;
it is caught only by the digest. Neither layer subsumes the other, which is why
the paper composes them rather than selecting one.

The honest gap is A3. BVI terminates a compromised credential and the
termination is immediate, terminal, and suppressible by neither the controller
nor the registrar acting alone. It does not retroactively invalidate sessions
anchored before compromise was detected. What it provides instead is an
immutable, consensus-timestamped record of exactly which sessions fell in that
window, which is what a dispute requires even though it is not prevention.

\subsection{Latency (RQ5)}

% PENDING: run scripts/measure-latency.js against Optimism Sepolia and insert
% p50/p95/p99 for the composite verification read and for anchor confirmation.
Verification latency and anchor confirmation latency are distinct quantities
and conflating them obscures the argument. The former is what must fit inside
the call-setup budget of requirement~R9: two \texttt{view} reads,
\texttt{resolve} and \texttt{isAuthorisedChannel}, performed by the recipient
over a public RPC with no account and no transaction. The latter is the time
from call teardown to inclusion of the anchor, and it sits off the call path
entirely --- which is precisely where BBCA's per-hop consensus does not.

\subsection{Throughput}

% Source: results/b4_throughput.json
At $91{,}535$~gas per anchor, an OP~Stack rollup at a $60$~million gas block
limit with two-second blocks sustains approximately $328$~anchors per second,
or $28.3$~million per day, for a chain carrying only BVI traffic. This is of
the same order as Australian daily call volume, so a single rollup at present
parameters does not clear national volume with comfortable margin. Batching
provides limited relief: amortising the $21{,}000$~gas intrinsic cost across a
batch reaches an asymptotic floor of $70{,}535$~gas per anchor, since the three
cold storage writes cannot be avoided, bounding the gain at $1.30\times$. A
raised block limit or partitioning across rollups by originating carrier are
the available levers. We record this as a deployment constraint.

\subsection{Threats to validity}

The digest figures are generated from synthetic speech and validate the
processing pipeline rather than establishing operational separability; they
require regeneration on a recorded narrowband corpus, and we expect the
separation to narrow, since synthesis produces cleaner spectro-temporal
gradients than speech does.
% PENDING: delete the sentence above once the real-corpus run is in.

The participation sweep is a simulation in which carriers attest independently
with equal probability. Real participation will be correlated by carrier size
and jurisdiction, which the model does not capture; we validated the simulator
against the closed-form expression and the two agree within $1.1$~percent, so
the discrepancy that matters is in the model's assumptions rather than in its
implementation.

Gas figures come from a local chain. They are reproducible and the toolchain is
pinned, but they are gas units rather than settled costs, and the fiat figures
quoted elsewhere depend on a gas price and an exchange rate that both move.

Finally, the revocation argument establishes that a revocation, once executed,
is immediate, terminal, observable, and suppressible by neither principal. It
does not establish prompt \emph{inclusion} on a rollup, which is a property of
the sequencer rather than of the contract, and is the reason status changes are
submitted to L1.
% PENDING: insert the measured forced-inclusion delay from the rollup's
% documentation. See docs/forced-inclusion.md.
